% 图 2-14：通路分层，层标签置于框外；连线走独立空白通道。
\begin{tikzpicture}[font=\footnotesize,
 b/.style={draw=softline,rounded corners=2pt,align=center,text width=3cm,inner sep=5pt,minimum height=1.2cm,fill=macblue},
 ar/.style={-{Stealth[length=1.4mm]},thick,draw=softline}]
\node[font=\small\bfseries,anchor=west] at (-7.25,.9) {主机与 PS：验证、调用与交互};
\node[b,fill=macpeach] (v1) at (-5.6,0) {参考实现 / KAT\\向量与期望值};
\node[b,fill=macpeach] (v2) at (-1.85,0) {run\_tests.py\\仿真与回归};
\node[b,fill=macpeach] (v3) at (1.9,0) {签核检查\\时序 / 资源 / IO};
\node[b,fill=macpeach] (v4) at (5.65,0) {AsconHost\\Jupyter 调用};
\node[b,text width=14.2cm,minimum height=1cm,fill=maclav] (protocol) at (0,-2) {当前链路：AXI-Lite CONTROL / MAILBOX $\rightarrow$ UART 字节 FIFO 帧协议};
\foreach \n in {v1,v2,v3,v4}{\draw[ar] (\n.south)--(protocol.north -| \n.south);}
\node[font=\small\bfseries,anchor=west] at (-7.25,-3.25) {RV32 裸机固件};
\node[b,fill=macgreen] (f1) at (-5.6,-4.2) {历史 ascon\_rv32\\纯软件};
\node[b,fill=macgreen] (f2) at (-1.85,-4.2) {历史\\\mbox{ascon\_hw\_rv32}\\MMIO 驱动};
\node[b,fill=macgreen] (f3) at (1.9,-4.2) {历史 ascon\_ise\_rv32\\定制指令};
\node[b,fill=macamber] (svc) at (5.65,-4.2) {\textbf{当前}\\\textbf{\mbox{pynq\_svc\_puf}}\\SW / MMIO / ISE\\+ PUF};
\draw[ar] (protocol.south -| svc.north)--(svc.north);
\node[font=\small\bfseries,anchor=west] at (-7.25,-5.65) {PL 硬件与计时};
\node[b,fill=macblue] (h1) at (-4.8,-6.65) {RV32 CPU\\csr\_file.mcycle};
\node[b,fill=macamber] (h2) at (0,-6.65) {Ascon MMIO / ISE\\两条硬件加速路径};
\node[b,fill=macblue] (h3) at (4.8,-6.65) {UART / CPI\\通信与事件观测};
% 从服务固件框右侧引出，沿右侧空白通道下行再绕到 CPU 底部，不经过任何框
\draw[ar] (svc.east)--(7.65,0 |- svc.east)--(7.65,-7.8)--(-4.8,-7.8)--(h1.south);
\draw[ar] (h1.east)--(h2.west);\draw[ar] (h2.east)--(h3.west);
\node[align=left,text width=14.2cm,font=\scriptsize,anchor=north west] at (-7.25,-8.35) {历史三套镜像与当前服务固件分列；当前性能计时来自 mcycle，CPI 事件计数不能替代总周期。\\上层工具的验证关系为软件工作流，不表示各脚本都直接访问 AXI-Lite。};
\end{tikzpicture}
